% Generated draft: requires mathematical and source review.
\documentclass[11pt]{article}
\usepackage[margin=25mm]{geometry}
\usepackage{fontspec}
\setmainfont{Latin Modern Roman}
\usepackage{amsmath,amssymb,mathtools,unicode-math}
\setmathfont{Latin Modern Math}
\usepackage{enumitem,needspace,xurl,hyperref}
\setlength{\parindent}{0pt}
\setlength{\parskip}{5pt}
\setlength{\emergencystretch}{3em}
\newcommand{\N}{\mathbb{N}}
\newcommand{\Z}{\mathbb{Z}}
\newcommand{\Q}{\mathbb{Q}}
\newcommand{\R}{\mathbb{R}}
\newcommand{\C}{\mathbb{C}}
\newcommand{\F}{\mathbb{F}}
\newcommand{\A}{\mathbb{A}}
\newcommand{\PP}{\mathbb{P}}
\newcommand{\E}{\mathbb{E}}
\newcommand{\eps}{\varepsilon}
\newcommand{\dd}{\,\mathrm{d}}
\begin{document}
\section*{UnsolvedMath ID: 30004413}
\textbf{Source problem code:} catalogue entry
\subsection{Definitions and mathematical statement}
Let \(T\) be a once-punctured torus and let \(G=\pi_1(T)\). Write \(\widehat{\Q}=\Q\cup\{\infty\}\) and \(\operatorname{PSL}(2,\C)=\operatorname{SL}(2,\C)/\{\pm I\}\). A representation \(\rho:G\to\operatorname{PSL}(2,\C)\) is discrete if its image is discrete in the usual topology, and nonfaithful if \(\ker\rho\ne\{1\}\). For a subset \(S\subseteq G\), its normal closure is
\[
\langle\!\langle S\rangle\!\rangle_G
=\bigcap_{\substack{N\trianglelefteq G\\S\subseteq N}}N.
\]
Under the usual slope interpretation, a marking of \(T\) indexes the unoriented isotopy classes of essential nonperipheral simple closed curves by \(r\in\widehat{\Q}\); \(\beta_r\in G\) denotes a representative of the corresponding conjugacy class, up to inversion. Changing the representative by conjugation or inversion does not affect the normal closures below. The supplied record does not explicitly define \(\beta_r\), so this interpretation requires confirmation against the cited report.

The record also does not define ``type-preserving.'' A common peripheral condition is that a loop around the puncture has nontrivial parabolic image. Whether the source requires only this condition or also restricts which other elements may have parabolic image is unresolved here. The conjecture below retains the source's type-preserving hypothesis without choosing between these interpretations.

\textbf{Conjectural statement, subject to these terminology qualifications.} For every nonfaithful, discrete, type-preserving representation
\[
\rho:G\longrightarrow\operatorname{PSL}(2,\C),
\]
one of the following holds:
\[
\exists r\in\widehat{\Q}\ \exists m\in\Z,\quad
m\ge2\quad\text{and}\quad
\ker\rho=\langle\!\langle\beta_r^m\rangle\!\rangle_G;
\]
or
\[
\begin{gathered}
\exists r_1,r_2\in\widehat{\Q}\ \exists m_1,m_2\in\Z,\qquad
r_1\ne r_2,\quad m_1,m_2\ge2,\\
\ker\rho=
\langle\!\langle\beta_{r_1}^{m_1},\beta_{r_2}^{m_2}\rangle\!\rangle_G.
\end{gathered}
\]
The alternatives are not asserted to be mutually exclusive. No uniqueness of slopes or exponents, or assertion that the exponents are exact orders of the corresponding images, is included in the supplied statement.
\subsection{Short English statement}
The supplied conjecture asks whether every discrete, noninjective, type-preserving representation of a once-punctured torus group into \(\operatorname{PSL}(2,\C)\) has kernel normally generated by a power of one simple loop, or by powers of two simple loops of distinct slopes, with all exponents at least \(2\). The meanings of the slope notation and type-preserving condition require source confirmation. The supplied open status is a dated curation label, not an independently verified conclusion.
\subsection{Sources}
\begin{itemize}
\item Local registry: supplied UnsolvedMath source record, permanent ID \(30004413\), source problem code \texttt{catalogue entry}, in the Oberwolfach Reports Open Problems collection. The project identifies \texttt{data/unsolvedmath/problems.json} and the supplemental \texttt{catalog/problems.json} as local registries; no file lookup was performed for this draft.
\item Supplied report citation: \emph{Low-dimensional Topology} (2021), \emph{Oberwolfach Reports}, pp.~465--516, DOI \texttt{10.4171/owr/2020/8}. \url{https://doi.org/10.4171/owr/2020/8}. The record gives proposed year 2020 and supplies no individual proposer. Its background describes Oberwolfach Reports as documenting mathematical workshops and their open-problem discussions.
\item Supplied literature source: Lee and Sakuma, ``Heckoid groups.'' \url{https://arxiv.org/abs/1206.4258}.
\item Supplied literature source: Lee and Sakuma, ``simple loops and Heckoid groups.'' \url{https://arxiv.org/abs/1402.6873}.
\item Supplied literature source: ``Oberwolfach Report 8/2020.'' \url{https://doi.org/10.4171/owr/2020/8}.
\item Supplied literature assessment, checked \textbf{2026-08-22}: the record labels the problem open and states that the kernel classification remains conjectural, with important Heckoid and simple-loop special cases proved. This is supplied dated curation, described in the record as literature triage derived from the supplied OWR curation; neither its status nor its claims have been independently verified here.
\item Supplied statement-verification metadata labels the statement \texttt{corrected\_verified} and reports that the official report was checked and that the punctured-torus setting and both conjectural kernel forms were retained. This records the supplied verification claim; it does not describe verification performed for this draft.
\end{itemize}
\end{document}
